% ============================================================
% Sección 2
% ============================================================
\section{2 EDO Lineal Homogénea}

\begin{tcolorbox}[enunciado]
\textbf{Clases:}
\begin{itemize}[leftmargin=*]
    \item Video 1) 2da parte Homogénea - Separable, Def. Edo Lineal, Homogénea y Ejemplo. 21 Sep26
    \item Video 2) EDO Lineal 1er Orden Homogénea y Variación de Parámetros MCA4 23Sep26
\end{itemize}
\end{tcolorbox}

\subsection*{ACTIVIDADES}

\begin{enumerate}[label=\arabic*., leftmargin=*]
    \item \textbf{Definición de EDO Lineal de 1er Orden} \\
    Pon la definición de EDO Lineal de primer orden dada en la clase del 21 de Septiembre (Video 1).
    \subsection*{Solución}
    % Escribe tu respuesta aquí

    \item ¿Qué es la Vida Media? (Video 1)
    \subsection*{Solución}
    % Escribe tu respuesta aquí

    \item \textbf{Descarga de una cola (de impresión o tareas) con tasa variable} \\
    Replica la información (la puedes resumir sin que pierda sentido) y solución que se dio a este problema en la clase del 23 de Septiembre (Video 2).
    \subsection*{Solución}
    % Escribe tu respuesta aquí

    \item \textbf{Pérdida de datos en una memoria volátil}

    Las memorias volátiles (como SDRAM) pierden su contenido cuando se interrumpe la alimentación eléctrica; incluso sin apagado total, los datos pueden degradarse con el tiempo debido a fugas de carga. Si se supone que la tasa de pérdida de datos es proporcional a la cantidad almacenada, se obtiene una EDO lineal homogénea. Este modelo es útil para determinar el tiempo máximo que se puede mantener la información sin refresco.

    La ecuación diferencial modelo es:
    \[ \frac{dD}{dt} + \delta D = 0 \]
    donde $D(t)$ es la cantidad de datos (bits) y $\delta > 0$ es la tasa de pérdida en unidades de $\text{s}^{-1}$.

    Si tenemos al inicio (condición inicial al tiempo $t=0$) una cantidad de datos $D(0) = D_{0}$, responde lo siguiente:

    \begin{enumerate}[label=(\alph*), leftmargin=*]
        \item Identifique $a(t)$.
        \subsection*{Solución}
        % Escribe tu respuesta aquí

        \item Halle $D(t)$.
        \subsection*{Solución}
        % Escribe tu respuesta aquí

        \item ¿En qué tiempo pierde el $99\%$ de los datos?
        \subsection*{Solución}
        % Escribe tu respuesta aquí
    \end{enumerate}
\end{enumerate}